\subsection{赋值域上的赋范代数}

定义 9. 若 A 是非离散赋值域 K 上的代数，则 A 上的范数 \(p(x)\) （把 A 看作 K 上的向量空间）称为与 A 的代数结构相容的，如果它定义的拓扑与 A 的环结构相容。K 上的代数赋有由与代数结构相容的范数所定义的结构后，称为赋范代数。

若 A 是 K 上的赋范代数，\(p(\boldsymbol{x})\) 是 A 上的范数，则双线性映射 \((\boldsymbol{x}, \boldsymbol{y}) \to \boldsymbol{x}\boldsymbol{y}\) （A × A 到 A 中的）由假设是连续的；于是由第 5 号的定理 1，存在实数 a \textgreater{} 0 使 \(p(\boldsymbol{x}\boldsymbol{y}) \leqslant a.p(\boldsymbol{x})p(\boldsymbol{y})\) 。把 \(p(\boldsymbol{x})\) 换成等价范数 \(a.p(\boldsymbol{x})\) ，因此我们总可设赋范代数 A 上的范数 \(||\boldsymbol{x}||\) 满足

\[
\left| \left| x y \right| \right| \leqslant \left| \left| x \right| \right|. \left| \left| y \right| \right|.\tag{5}
\]

由 (5) 用归纳法推出，对每个整数 \(n > 0\) 有

\[
\left| \left| x ^ {n} \right| \right| \leqslant \left| \left| x \right| \right| ^ {n}.\tag{6}
\]

例。1) 设 K 是赋值除环，\(K'\) 是 K 的中心的一个子域，使得在 \(K'\) 上 K 的绝对值 \(|x|\) 的迹不是 \(K'\) 上的平凡绝对值。于是 K 以 \(|x|\) 为范数时是 \(K'\) 上的赋范代数。

\begin{enumerate}
\def\labelenumi{\arabic{enumi})}
\setcounter{enumi}{1}
\item
  设 K 是非离散赋值域，\(\mathbf{M}_{n}(\mathbf{K})\) 是 K 上 n 阶方阵的环。把 \(\mathbf{M}_{n}(\mathbf{K})\) 看作 K 上的向量空间时，它同构于 \(K^{n^{2}}\) 。若对每个 \(X = (x_{ij}) \in \mathbf{M}_{n}(\mathbf{K})\) 定义 \(\|X\| = \sup_{i,j} |x_{ij}|\) ，则 \(\|X\|\) 是 \(\mathbf{M}_{n}(\mathbf{K})\) 上的范数，且由这个范数定义的拓扑是 \(K^{n^{2}}\) 上的积拓扑（命题 10）；由此（由于 K 上任意多个变量的多项式的连续性）推出这个范数与 \(\mathbf{M}_{n}(\mathbf{K})\) 的 K-代数结构相容。
\item
  集 \(\mathcal{B}(\mathbf{X};\mathbf{K})\) ——由全体函数 \(f\) （在集 \(\mathbf{X}\) 上、取值于非离散赋值域 \(\mathbf{K}\) ，且使 \(x\to |f(x)|\) 在 \(\mathbf{X}\) 上有界）组成——是 \(\mathbf{K}\) 上的代数；范数 \(\| f\| = \sup_{x\in \mathbf{X}}|f(x)|\) 与 \(\mathcal{B}(\mathbf{X};\mathbf{K})\) 的环结构相容，因为我们有 \(\| fg\| \leqslant \| f\| .\| g\|\) （参见第 X 章，§ 1）。
\end{enumerate}

设 \(\mathfrak{a}\) 是赋范代数 \(\mathbf{A}\) 中的闭双边理想。若在商代数 \(\mathbf{A} / \mathfrak{a}\) 中令 \(||\dot{\mathbf{x}}|| = \inf_{\mathbf{x}\in \dot{\mathbf{x}}}||\mathbf{x}||\) ，就得到 \(\mathbf{A} / \mathfrak{a}\) 上的一个范数，它定义的拓扑是 \(\mathfrak{a}\) 对 \(\mathbf{A}\) 的拓扑所作的商（命题 9）；由于这个商拓扑与 \(\mathbf{A} / \mathfrak{a}\) 的商环结构相容（第 III 章，§ 6，第 4 号），由此推出商代数 \(\mathbf{A} / \mathfrak{a}\) （赋有范数 \(||\mathbf{x}||\) ）是赋范代数。

同样，若 \((\mathbf{A}_i)_{1 \leqslant i \leqslant n}\) 是 \(n\) 个赋值域 \(\mathbf{K}\) 上的赋范代数组成的族，且在积代数 \(\mathbf{A} = \prod_{i=1}^{n} \mathbf{A}_i\) 中令 \(||\boldsymbol{x}|| = \sup_{i} ||\boldsymbol{x}_i||\) （其中 \(\boldsymbol{x} = (\boldsymbol{x}_i)\) ），则得到 \(\mathbf{A}\) 上的一个范数，它定义诸 \(\mathbf{A}_i\) 的拓扑之积（命题 10）；由于这个拓扑与 \(\mathbf{A}\) 的环结构相容（第 III 章，§ 6，第 4 号），由此推出积代数 \(\mathbf{A}\) （赋有范数 \(||\boldsymbol{x}||\) ）是赋范代数。

设 A 是赋值域 K 上的赋范代数。A 的完备化 \(\hat{\mathbf{A}}\) （第 III 章，§ 6，第 5 号，命题 6）还赋有 \(\hat{\mathbf{K}}\) 上的向量空间结构（第 3 号，命题 8），并且由恒等式延拓原理显然有 \(t(xy) = (tx)y = x(ty)\) （对一切 \(t \in \hat{\mathbf{K}}\) 与一切 \(x, y \in \hat{\mathbf{A}}\) ）。因此 \(\hat{\mathbf{A}}\) 是 \(\hat{\mathbf{K}}\) 上的代数；另一方面（第 3 号，命题 8），A 上的范数由连续性扩张为定义 \(\hat{\mathbf{A}}\) 拓扑的范数，因此 \(\hat{\mathbf{A}}\) 赋有这个范数后是域 \(\hat{\mathbf{K}}\) 上的赋范代数。

若 \((\pmb{x}_{\lambda})_{\lambda \in \mathbf{L}}\) 与 \((\pmb{y}_{\mu})_{\mu \in \mathbf{M}}\) 是赋范代数 \(\mathbf{A}\) 中两个绝对可和族，则族 \((\pmb{x}_{\lambda}\pmb{y}_{\mu})_{(\lambda, \mu) \in \mathbf{L} \times \mathbf{M}}\) 绝对可和，因为 \(||\pmb{x}_{\lambda}\pmb{y}_{\mu}|| \leqslant ||\pmb{x}_{\lambda}|| \cdot ||\pmb{y}_{\mu}||\) （第 IV 章，§ 7，第 3 号，命题 1）；若此外 \(\mathbf{A}\) 完备，则三个族都可和，且我们有

\[
\sum_ {(\lambda , \mu) \in \mathbf {L} \times \mathbf {M}} x _ {\lambda} y _ {\mu} = \left(\sum_ {\lambda \in \mathbf {L}} x _ {\lambda}\right) \left(\sum_ {\mu \in \mathbf {M}} y _ {\mu}\right)
\]

这由左端和式的结合性得出（第 III 章，§ 5，第 3 号，公式 (2)）。

若赋范代数 A 有单位元 \(e \neq 0\) ，则映射 \(t \rightarrow te\) 是 K 的域结构到 A 的子域 Ke 的域结构上的同构；这个同构也是 K 的拓扑域结构到 Ke 的拓扑域结构上的同构（后者的拓扑由 A 的拓扑诱导），因为 A 的范数在 Ke 上的限制 \(\|te\|\) 是与绝对值

\[
| t | = \frac {\mathrm{I}}{| | e | |} | | t e | |.
\]

若 \(||e|| = 1\) ，则 \(||te|| = |t|\) ，于是可把赋值域 \(K\) 与赋范子域 \(Ke\) （\(A\) 的）等同，特别地，可把 \(A\) 的单位元记作符号 \(I\) 。

以下我们只讨论有单位元 e 且范数满足不等式 (5) 的赋范代数；在这个不等式中取 x = y = e ，便得 \(\|e\| \geqslant 1\) 。

命题 12. 若通项为 \(z^n\) 的级数在 A 中收敛，则 \(e - z\) 是 A 的可逆元，并且我们有

\[
(e - z) ^ {- 1} = \sum_ {n = 0} ^ {\infty} z ^ {n}.\tag{7}
\]

反之，若 \(||z|| < 1\) 且 \(e - z\) 是 \(A\) 中的可逆元，则通项为 \(z^n\) 的级数收敛，且公式 (7) 成立。

对每个 \(p > 0\) 我们有

\[
(e - z) \sum_ {n = 0} ^ {p} z ^ {n} = e - z ^ {p + 1}.\tag{8}
\]

若通项为 \(z^{n}\) 的级数收敛且 y 是它的和，则 \(z^{n}\) 当 \(n \to +\infty\) 时趋于 o；于是在 (8) 中取极限便有 \((e - z)y = e\) ；类似地可证 \(y(e - z) = e\) ，从而 \(y = (e - z)^{-1}\) （注意论证的这一部分在任何有单位元的拓扑环中都成立）。

反之，若 \(||z|| < 1\) ，则由于 \(||z^{p+1}|| \leqslant ||z||^{p+1}\) ，\(z^{p+1}\) 当 \(p \to +\infty\) 时趋于 o；用 \((e - z)^{-1}\) 左乘 (8) 的两边并令 \(p\) 趋于无穷，便知通项为 \(z^n\) 的级数收敛且以 \((e - z)^{-1}\) 为其和。

推论. 设 A 是完备赋范代数。则对每个 \(z \in \mathbf{A}\) （满足 \(||z|| < 1\) ），\(e - z\) 是 A 中的可逆元。

通项为 \(z^n\) 的级数绝对收敛，因为 \(||z^n|| \leqslant ||z||^n\) （对 \(n > 0\) ），从而收敛，因为 \(A\) 完备（第 6 号，命题 11）。

命题 13. 设 G 是完备赋范代数 A 的可逆元群。则 G 是 A 的开子集；A 的拓扑在 G 上诱导的拓扑与 G 的群结构相容；且赋有这个拓扑的 G 是完备群（关于它的两个一致结构中的每一个）。

命题 12 的推论表明 G 包含 e 在 A 中的一个邻域 V；于是对每个 \(x_{0} \in G\) ，\(x_{0}V\) 的元素都是可逆元，且 \(x_{0}V\) 是 \(x_{0}\) 在 A 中的邻域，因为 \(x \to x_{0}x\) 是 A 到自身的同胚（\(x_{0}\) 是 A 的可逆元）。因此 G 在 A 中是开的。

要证明 \(\mathbf{G}\) 上由 \(\mathbf{A}\) 的拓扑诱导的拓扑与 \(\mathbf{G}\) 的群结构相容，只需证明函数 \(x^{-1}\) 在 \(\mathbf{G}\) 上连续。设 \(x_0 \in \mathbf{G}\) ，对每个 \(x \in \mathbf{G}\) ，把 \(\boldsymbol{x}\) 写成形式 \(x = x_0 (e + u)\) ，于是 \(u = x_0^{-1}(x - x_0)\) ；则 \(||\boldsymbol{u}|| \leqslant ||x_0^{-1}|| \cdot ||\boldsymbol{x} - x_0||\) ，因而若 \(||\boldsymbol{x} - x_0|| \leqslant 1 / ||x_0^{-1}||\) ，就有 \(||\boldsymbol{u}|| \leqslant 1\) ，\(e + u = x_0^{-1}x\) 是可逆元，通项为 \((-1)^n u^n\) 的级数绝对收敛，且

\[
\boldsymbol {x} ^ {- 1} = (\boldsymbol {e} + \boldsymbol {u}) ^ {- 1} \boldsymbol {x} _ {0} ^ {- 1} = \boldsymbol {x} _ {0} ^ {- 1} + \left(\sum_ {n = 1} ^ {\infty} (- 1) ^ {n} \boldsymbol {u} ^ {n}\right) \boldsymbol {x} _ {0} ^ {- 1},\tag{9}
\]

由此推出

\[
\begin{array}{r l} & {| | x _ {0} ^ {- 1} - x _ {0} ^ {- 1} | | \leqslant \left\| \sum_ {n = 0} ^ {\infty} (- \mathrm{I}) ^ {n} u ^ {n} \right\|. | | u | |. | | x _ {0} ^ {- 1} | |} \\ & {\quad \leqslant \left\| \sum_ {n = 0} ^ {\infty} (- \mathrm{I}) ^ {n} u ^ {n} \right\|. | | x _ {0} ^ {- 1} | | ^ {2}. | | x - x _ {0} | |.} \end{array}
\]

当 \(x\) 趋于 \(x_0\) 时，\(||x - x_0||\) 趋于 \(o\) ，并且由于

\[
\left\| \sum_ {n = 0} ^ {\infty} (- 1) ^ {n} \boldsymbol {u} ^ {n} \right\| \leqslant | | \boldsymbol {e} | | + \frac {| | \boldsymbol {u} | |}{1 - | | \boldsymbol {u} | |}
\]

保持有界，由此得出 \(x^{-1}\) 趋于 \(x_{0}^{-1}\) 。

最后，为证明 G 的左一致结构是完备的，我们来证明关于这个一致结构的每个柯西滤子 F 都是关于 A 的加法一致结构的柯西滤子且收敛于 G 的一点。对每个 \(\varepsilon\) （满足 \(0 < \varepsilon < 1\) ），有集 \(\mathbf{M} \in \mathfrak{F}\) 使得 \(||x^{-1}y - e|| \leqslant \varepsilon\) 对 M 中一切 \(x, y\) 成立，即使得

\[
\left| \left| \mathbf {y} - \mathbf {x} \right| \right| \leqslant \varepsilon \left| \left| \mathbf {x} \right| \right|.
\]

设 \(\pmb{a}\) 是 M 的一点；对每个 \(x \in M\) 有 \(||x - a|| \leqslant \varepsilon ||a||\) ，因此 \(||x|| \leqslant (1 + \varepsilon)||a||\) 。另一方面，存在集 \(N \subset M\) （属于 \(\mathfrak{F}\) ），使得 \(||x^{-1}y - e|| \leqslant \varepsilon / (1 + \varepsilon)||a||\) 对 N 中一切 \(x\) 与 \(y\) 成立；由此得 \(||y - x|| \leqslant \varepsilon ||x|| / (1 + \varepsilon)||a|| \leqslant \varepsilon\) ，这表明 \(\mathfrak{F}\) 是关于 A 的加法一致结构的柯西滤子，因此由于 A 完备而收敛于一点 \(x_0\) 。由于 \(x_0\) 是 \(\mathfrak{F}\) 的极限，有 \(||x^{-1}x_0 - e|| \leqslant \varepsilon\) （对一切 \(x \in M\) ，由不等式延拓原理）；由 \(\varepsilon > 1\) 推出 \(x^{-1}x_0\) 是 A 中的可逆元；因此 \(x_0\) 是可逆元，即 \(x_0 \in G\) 。

命题 14. 在完备赋值除环中，非零元素所成的乘法群是完备群。

证明与命题 13 类似；只需把 A 的范数换成所考虑除环的绝对值。

注意我们不能直接应用命题 13，因为（非交换的）赋值除环不一定是非离散赋值域上的代数（绝对值在除环中心上的限制可能是平凡的）。

注。命题 13 对非完备的赋范代数未必成立。例如，在代数 \(\mathcal{C}(\mathbf{I};\mathbf{R})\) （由 \(\mathbf{I} = [\mathrm{o},\mathbf{i}]\) 上一切有限连续实值函数组成）中{[}范数取 \(\| x\| = \sup_{t\in \mathbf{I}}|x(t)|\) ，t 的一切多项式（限制在 I 上）所成的子代数 P 不完备；若 \(x(t)\) 是任一非常数多项式，则 \(1 + \varepsilon x\) 当 \(\varepsilon\) 任意小时任意接近 P 的单位元 i，但 \(1 + \varepsilon x\) 不是 P 中的可逆元。然而，若 A 是非完备赋范代数，G 是它的可逆元群，\(\hat{A}\) 是 A 的完备化，则 G 是 \(\hat{A}\) 的可逆元群的子群，因此 A 的拓扑在 G 上诱导的拓扑与 G 的群结构相容。

